\documentclass[twocolumn]{article}
\usepackage[margin=0.7in,columnsep=0.25in]{geometry}
\usepackage{eqannotate}
\begin{document}
\section*{Clustered terms}
\begin{annotatedequation}
F=\eqmark[blue]{a}{a}+\eqmark[red]{b}{b}+\eqmark[green]{c}{c}+\eqmark[orange]{d}{d}+\eqmark[purple]{e}{e}+\eqmark[cyan]{f}{f}
\eqannotate{a}{First long label}
\eqannotate{b}{Second label}
\eqannotate{c}{Third long explanatory label}
\eqannotate{d}{Fourth label}
\eqannotate{e}{Fifth explanatory label}
\eqannotate{f}{Sixth label}
\end{annotatedequation}

\section*{Geometry-driven sides}
\begin{annotatedequation}
q(x)=\frac{\eqmark[yellow]{num}{(x-\mu)^2}+\eqmark[cyan]{num2}{\alpha x}}{\eqmark[blue]{den}{2\sigma^2}+\eqmark[red]{den2}{\beta}}
\eqannotate{num}{Squared residual}
\eqannotate{num2}{Numerator term}
\eqannotate{den}{Variance scale}
\eqannotate{den2}{Denominator offset}
\end{annotatedequation}

\section*{Forced side hint}
\begin{annotatedequation}
G=\eqmark[green]{x}{x}+\eqmark[orange]{y}{y}
\eqannotate[prefer=below]{x}{Forced below}
\eqannotate[prefer=above]{y}{Forced above}
\end{annotatedequation}
\end{document}
